\chapter{REFERENCIAS}
\thispagestyle{empty}

% =====================================================================
%  PENDIENTE  --  Capitulo 3
%  Este capitulo es un guia de la plantilla y contiene un ejemplo de
%  como citar. Cuando el documento este completo:
%    1) La bibliografia final se genera automaticamente al final del
%       documento con \printbibliography{} (NO duplicarla aqui).
%    2) Este capitulo puede conservarse como "Guia de citacion" o
%       eliminarse comentando \include*{cap03} en ReportePI.tex.
%
%  Formato de referencia utilizado: APA (configurado en el preambulo
%  con \usepackage[natbib=true,backend=biber,sorting=nyt,style=apa]{biblatex}).
%  Para agregar una referencia nueva: editar referencias.bib y citarla
%  con \cite{clave}.
%
%  Ejemplo de entrada BibTeX:
%
%  @article{velasquez2017relation,
%    title={Relation between BMI and Exit Time of a Building in an Emergency Situation},
%    author={Velásquez, Washington and Munoz-Arcentales, Andres and Salvachúa, Joaquín},
%    journal={Engineering Letters},
%    volume={25},
%    number={4},
%    year={2017}
%  }
%
%  Cita en el texto: El articulo presentado por Velásquez, et al.
%  \cite{velasquez2017relation}.
% =====================================================================

\section{C\'omo citar en este documento}

Las referencias utilizadas en el texto se invocan con el comando \verb|\cite{clave}|,
donde \emph{clave} corresponde al identificador de la entrada en el archivo
\verb|referencias.bib|. La bibliograf\'ia final se genera autom\'aticamente al
final del documento con el comando \verb|\printbibliography{}|, empleando el
estilo APA definido en el pre\'ambulo.

Las fuentes consultadas para la construcci\'on del marco te\'orico incluyen,
entre otras, la especificaci\'on de las capas de las redes programables
\cite{rfc7426}, los trabajos originales sobre el controlador de red empleado
\cite{ryu2015onosr,openvswitch2014}, la literatura sobre aprendizaje autom\'atico
aplicado a la detecci\'on de intrusiones \cite{lightgbm2017} y las gu\'ias
institucionales de gesti\'on de incidentes de seguridad \cite{nist80061r2,nistcsf20}.
